\documentclass[12pt]{article}
\usepackage[margin=1in]{geometry}
\usepackage{amsmath,amssymb,amsthm}
\usepackage{hyperref}

\newtheorem{theorem}{Theorem}
\newtheorem{lemma}[theorem]{Lemma}
\newtheorem{proposition}[theorem]{Proposition}
\newtheorem{corollary}[theorem]{Corollary}
\theoremstyle{definition}
\newtheorem{definition}[theorem]{Definition}
\theoremstyle{remark}
\newtheorem{remark}[theorem]{Remark}

\newcommand{\R}{\mathbb{R}}
\newcommand{\C}{\mathbb{C}}
\newcommand{\Z}{\mathbb{Z}}
\newcommand{\E}{\mathbb{E}}
\newcommand{\Var}{\mathrm{Var}}
\newcommand{\Tr}{\mathrm{Tr}}
\newcommand{\ph}{\varphi}
\newcommand{\La}{\Lambda}
\newcommand{\im}{\mathrm{Im}}
\newcommand{\re}{\mathrm{Re}}

\title{Complex Actions, Sign Problems, and Contour Deformations in QFT}
\author{M.~R.~Douglas}
\date{April 2026}

\begin{document}
\maketitle

\begin{abstract}
We survey the sign problem in Euclidean lattice QFT and the
contour deformation methods used to address it, with emphasis on
the Lefschetz thimble decomposition and its application to the
O($N$) model at large~$N$. We connect the steepest descent contour
strategy of \texttt{mass-gap-v3} to the broader thimble literature,
identify the key mathematical steps needed for a rigorous proof,
and discuss both positive results and fundamental obstructions.
\end{abstract}

\tableofcontents

%% ============================================================
\section{The sign problem in Euclidean QFT}
%% ============================================================

\subsection{General setup}

A lattice QFT partition function has the form
\begin{equation}\label{eq:Z}
  Z = \int_{\R^n} e^{-S(\phi)}\,d\phi,
\end{equation}
where $\phi \in \R^n$ is the field configuration ($n = N|\La|$ for
$N$-component fields on lattice~$\La$) and $S : \R^n \to \R$ is the
Euclidean action. When $S$ is real, $e^{-S} > 0$ and the integrand
defines a probability measure (after normalization). Observables are
expectations under this measure.

The \textbf{sign problem} arises when a change of variables or
auxiliary field introduces a \emph{complex} action $S : \R^n \to \C$,
so $e^{-S}$ is complex-valued. The integrand is no longer a positive
measure, and Monte Carlo methods fail because importance sampling
requires positive weights.

\subsection{Sources of complex actions}

Complex actions arise in several settings:
\begin{itemize}
\item \textbf{Finite density}: Chemical potential $\mu$ gives
  $\det(D+\mu\gamma_0)$ which is complex for QCD.
\item \textbf{Real-time evolution}: The Minkowski action $iS_M$ is
  purely imaginary.
\item \textbf{Hubbard-Stratonovich transformation}: The HS identity
  for the Euclidean quartic $e^{-\lambda a^2}$ requires imaginary
  coupling $e^{i\sigma a}$ (Section~\ref{sec:HS}).
\item \textbf{Topological terms}: The $\theta$-term $e^{i\theta Q}$
  is complex for $\theta \ne 0$.
\item \textbf{Fermion determinant}: $\det(D)$ can be complex for
  chiral or finite-density theories.
\end{itemize}

\subsection{Why contour deformation?}

Since $S(\phi)$ extends to a holomorphic (or meromorphic) function
$S : \C^n \to \C$, the integral~\eqref{eq:Z} can be deformed from
$\R^n$ to any $n$-dimensional cycle $\Gamma \subset \C^n$ in the
same homology class, provided $\Gamma$ avoids singularities and the
integrand decays at infinity:
\begin{equation}
  Z = \int_\Gamma e^{-S(\sigma)}\,d^n\sigma.
\end{equation}
The goal is to choose $\Gamma$ so that $e^{-S}$ is as close to
positive (non-oscillatory) as possible on~$\Gamma$.

%% ============================================================
\section{The HS sign problem}\label{sec:HS}
%% ============================================================

\subsection{The Euclidean quartic requires imaginary coupling}

For the O($N$) LSM, the quartic interaction enters as $e^{-\lambda a^2}$
with $a = |\ph|^2/N - v^2$. The HS identity:
\begin{equation}
  e^{-\lambda a^2}
    = \frac{1}{\sqrt{4\pi\lambda}}\int_{-\infty}^{\infty}
      e^{-\sigma^2/(4\lambda)\;+\;i\sigma a}\,d\sigma.
\end{equation}
The coupling $i\sigma a$ is \textbf{imaginary} because $e^{-\lambda a^2}$
is a Gaussian in $a$ (not $e^{+\lambda a^2}$). This is forced by
Euclidean signature.

\subsection{The $\sigma$-effective action}

After applying HS at each site and integrating out $\ph$:
\begin{equation}\label{eq:f}
  Z = c\int_{\R^{|\La|}} e^{f(\sigma)}\,d\sigma, \qquad
  f(\sigma) = -\sum_x\frac{\sigma(x)^2}{4\lambda}
    - \frac{N}{2}\Tr\log(-\Delta+2i\sigma z)
    - iv^2\sum_x\sigma(x),
\end{equation}
where $z = 1/\sqrt{N}$. This integral converges on $\R^{|\La|}$ (the
Gaussian factor provides decay, and $|\det(-\Delta+2i\sigma z)^{-1}|
\le \det(-\Delta)^{-1}$), but the integrand is \textbf{complex}.

\subsection{Comparison: NLSM vs LSM}

For the NLSM (nonlinear sigma model with constraint $|\ph|^2 = N\beta$),
a real mass $m_0^2$ can be added to the GFF for free (it is constant
on the constraint surface). The operator $-\Delta + m_0^2 - 2i\alpha z$
has positive real part, and the FK bound gives single-$\sigma$ decay
$|G_\sigma(x,0)| \le G_{m_0}(x,0) \sim e^{-m_0|x|}$.

For the LSM: the add-subtract trick cancels the real mass, leaving
$-\Delta - 2i\sigma$ with \textbf{no real part} beyond the kinetic
term. The mass gap comes entirely from the oscillatory
$\sigma$-average. See \texttt{mass-gap-v3.tex}, Sections~3--4.

%% ============================================================
\section{Lefschetz thimbles}
%% ============================================================

\subsection{Morse theory for complex integrals}

Consider a holomorphic function $f : \C^n \to \C$ (the exponent
in the partition function). The critical points of $f$ are the
solutions of $\partial f/\partial\sigma_i = 0$. For each critical
point $\sigma_*$, the \textbf{Lefschetz thimble} $\mathcal{J}_{\sigma_*}$
is the union of gradient flow lines of $\re\,f$ that flow \emph{into}
$\sigma_*$ as $t \to +\infty$:
\begin{equation}
  \mathcal{J}_{\sigma_*} = \Bigl\{\sigma(0) :
    \frac{d\sigma}{dt} = \overline{\nabla f(\sigma)},\;
    \lim_{t\to+\infty}\sigma(t) = \sigma_*\Bigr\}.
\end{equation}

Key properties of $\mathcal{J}_{\sigma_*}$:
\begin{enumerate}
\item $\mathcal{J}_{\sigma_*}$ is an $n$-dimensional (real)
  submanifold of~$\C^n$.
\item $\im\,f = \im\,f(\sigma_*)$ is \textbf{constant} on
  $\mathcal{J}_{\sigma_*}$ (the flow preserves the imaginary part).
\item $\re\,f$ is \textbf{maximized} at $\sigma_*$ on
  $\mathcal{J}_{\sigma_*}$ (steepest descent from the saddle).
\item The integrand $e^{f(\sigma)}$ is proportional to a
  \textbf{real positive} function times the constant phase
  $e^{i\,\im f(\sigma_*)}$ on~$\mathcal{J}_{\sigma_*}$.
\end{enumerate}

\subsection{Thimble decomposition}

The original integration cycle $\R^n$ can be decomposed:
\begin{equation}\label{eq:thimble_decomp}
  [\R^n] = \sum_{\sigma_*} n_{\sigma_*}\,[\mathcal{J}_{\sigma_*}],
\end{equation}
where $n_{\sigma_*} \in \Z$ are \textbf{intersection numbers}
(the intersection of $\R^n$ with the \emph{dual} thimble
$\mathcal{K}_{\sigma_*}$, defined by the ascending gradient flow).
Therefore:
\begin{equation}
  Z = \sum_{\sigma_*} n_{\sigma_*} \int_{\mathcal{J}_{\sigma_*}}
    e^{f(\sigma)}\,d\sigma
  = \sum_{\sigma_*} n_{\sigma_*}\,e^{i\,\im f(\sigma_*)}\,
    Z_{\sigma_*},
\end{equation}
where $Z_{\sigma_*} = \int_{\mathcal{J}_{\sigma_*}}
e^{\re f(\sigma) - \re f(\sigma_*)}\,|d\sigma| > 0$ is
\textbf{real and positive}.

\subsection{When is the sign problem solved?}

The sign problem is eliminated on a single thimble (since the integrand
is proportional to a positive function). But the \emph{inter-thimble}
sign problem remains if multiple thimbles contribute with different
phases $e^{i\,\im f(\sigma_*)}$.

The sign problem is fully solved when:
\begin{enumerate}
\item[(a)] Only one thimble has $n_{\sigma_*} \ne 0$, or
\item[(b)] All contributing thimbles have the same phase
  $\im f(\sigma_*)$, or
\item[(c)] One thimble dominates exponentially:
  $\re f(\sigma_*) \gg \re f(\sigma'_*)$ for all other
  contributing critical points~$\sigma'_*$.
\end{enumerate}

\subsection{Stokes phenomenon}

As parameters (coupling constants, volume) vary, the thimble
decomposition~\eqref{eq:thimble_decomp} can jump: when two critical
points have $\im f(\sigma_*) = \im f(\sigma'_*)$, a gradient flow
line connects them, and the intersection numbers change. These
\textbf{Stokes walls} are codimension-1 surfaces in parameter space.

Cao-Santachiara-Usciati~\cite{CSU23} found infinitely many Stokes
walls for lattice Liouville theory as $c \to (-\infty,1]$, showing
the thimble structure can be very intricate.

%% ============================================================
\section{Contour deformation methods in lattice QFT}
%% ============================================================

\subsection{Exact thimble methods}

The pioneering work of Cristoforetti-Scorzato et al.\ proposed
Monte Carlo sampling directly on a Lefschetz thimble. The main
challenge: computing the \textbf{residual phase} (the Jacobian
of the map from $\R^n$ to $\mathcal{J}_{\sigma_*}$), which is
expensive ($O(n^3)$ per configuration).

\subsection{Generalized thimble / gradient flow}

Instead of the exact thimble, one can use the \textbf{flowed
manifold}:
\begin{equation}
  \mathcal{M}_T = \bigl\{\sigma(T) : \sigma(0) \in \R^n,\;
    \dot\sigma = \overline{\nabla f(\sigma)}\bigr\},
\end{equation}
obtained by flowing the real domain along the anti-holomorphic
gradient of $f$ for time~$T$. As $T \to \infty$:
$\mathcal{M}_T \to \bigcup_{\sigma_*} \mathcal{J}_{\sigma_*}$.
For finite~$T$: $\mathcal{M}_T$ is a smooth deformation of $\R^n$
that already reduces the sign problem.

\textbf{Advantage}: No need to identify critical points or compute
intersection numbers. The flow is a well-posed ODE.

\subsection{Worldvolume HMC (Fukuma et al.)}

Fukuma and collaborators~\cite{Fukuma20,Fukuma23,Fukuma25,Fukuma26}
developed the most systematic computational approach:
\begin{enumerate}
\item Foliate $\C^n$ by a continuum family of integration surfaces
  $\{\mathcal{M}_T\}_{T \ge 0}$.
\item Introduce a symplectic structure on the tangent bundle of
  each~$\mathcal{M}_T$.
\item Perform HMC (Hamiltonian Monte Carlo) on $\mathcal{M}_T$,
  treating the surface as a Riemannian manifold.
\end{enumerate}
No Jacobian computation is needed for configuration generation
(only for reweighting). Extended to group manifolds (gauge theories)
in~\cite{Fukuma26}. Tested on the Hubbard model~\cite{Fukuma25b}
with severe sign problem.

\subsection{Constant contour shifts}

The simplest contour deformation: $\sigma \to \sigma + i\delta$
(constant imaginary shift). Ostmeyer et al.~\cite{Ostmeyer23}
proved this completely eliminates the sign problem in the Hubbard
model at large chemical potential. Gaentgen et al.~\cite{Gaentgen24}
applied this to C$_{20}$ and C$_{60}$ fullerenes.

\subsection{Fundamental limitations}

Lawrence-Yamauchi~\cite{LY24} proved that for 2D Abelian lattice
gauge theory, the sign problem is \textbf{exponential} and
\textbf{cannot be removed by any contour deformation}. The
obstruction is topological: the compact gauge group $U(1)$ has
$\pi_1(U(1)) = \Z$, creating nontrivial winding sectors that
contribute with different phases.

\begin{remark}
This negative result does \textbf{not} apply to scalar field theories
like the O($N$) model. The field space $\R^n$ is contractible
($\pi_k(\R^n) = 0$ for all~$k$), so there is no topological
obstruction to contour deformation. The question for scalar theories
is analytic (singularity avoidance, decay at infinity), not topological.
\end{remark}

%% ============================================================
\section{Application to the O($N$) model at large $N$}
%% ============================================================

\subsection{The exponent and its critical points}

The exponent~\eqref{eq:f} extends to a holomorphic function
$f : \C^{|\La|} \to \C$ (away from the singularities of the
log~det). The critical point equation $\partial f/\partial\sigma(x) = 0$
gives:
\begin{equation}\label{eq:gap}
  \frac{\sigma_*(x)}{2\lambda}
    = -iNz\,(-\Delta+2i\sigma_* z)^{-1}(x,x) - iv^2.
\end{equation}
For the translation-invariant saddle $\sigma_* = \sigma_0$
(constant), this becomes the \textbf{gap equation}:
\begin{equation}
  \frac{\sigma_0}{2\lambda} = -iz\bigl[N\,C_{\sigma_0}(x,x) + v^2\bigr],
  \qquad C_{\sigma_0} = (-\Delta+2i\sigma_0 z)^{-1}.
\end{equation}
At $\sigma_0 = 0$ this reads $0 = -iz[NC_0(0,0)+v^2]$, which
determines $v^2$ (or equivalently, the bare coupling).

\subsection{Single-thimble dominance at large $N$}

The exponent scales as $f \sim N\cdot g(\sigma)$ for a function~$g$
independent of~$N$. Therefore:
\begin{itemize}
\item The Hessian scales as $f'' \sim N$, giving saddle-point
  width $\sim 1/\sqrt{N}$.
\item At competing saddle points $\sigma'_*$: the action difference
  $\re f(\sigma_*) - \re f(\sigma'_*) \sim N\cdot\Delta g > 0$ grows
  linearly in~$N$.
\item By condition~(c) of Section~3.3: the dominant thimble
  contribution exceeds all others by a factor $e^{cN}$ for some $c > 0$.
\end{itemize}

\textbf{Conclusion}: For $N$ sufficiently large, the single thimble
$\mathcal{J}_{\sigma_*}$ through the dominant saddle captures $Z$
up to exponentially small corrections:
\begin{equation}
  Z = n_{\sigma_*}\,e^{i\,\im f(\sigma_*)}\,Z_{\sigma_*}
    \bigl(1 + O(e^{-cN})\bigr).
\end{equation}

\subsection{Properties of the dominant thimble}

On $\mathcal{J}_{\sigma_*}$:
\begin{enumerate}
\item The integrand $e^{f(\sigma)}$ is proportional to a
  \textbf{positive} function (Section~3.1, property~4).
\item Near $\sigma_*$: $f \approx f(\sigma_*) - \tfrac{1}{2}
  (\delta\sigma,\,H\,\delta\sigma)$ with $H = -f''(\sigma_*) > 0$
  (positive definite Hessian). The measure is \textbf{log-concave}
  near the saddle.
\item The $\ph$-coupling becomes effectively \textbf{real} on the
  thimble (since $\im f = \text{const}$ constrains the imaginary
  parts of $\sigma$). This is needed for the FK bound.
\end{enumerate}

\subsection{The proof strategy}

Combining the thimble structure with the Brascamp-Lieb and
Feynman-Kac ingredients:

\begin{enumerate}
\item \textbf{Thimble existence}: Show $\mathcal{J}_{\sigma_*}$ is a
  smooth $|\La|$-dimensional manifold avoiding the singularities
  $\det(-\Delta+2i\sigma z) = 0$.
\item \textbf{Single-thimble dominance}: For $N > N_0$, the
  intersection number $n_{\sigma_*} = 1$ and all other thimbles
  are suppressed by $e^{-cN}$.
\item \textbf{Cauchy's theorem}: The deformation from $\R^{|\La|}$
  to $\mathcal{J}_{\sigma_*}$ is valid (no singularities crossed,
  boundary terms vanish).
\item \textbf{Brascamp-Lieb on the thimble}: The measure
  $e^{\re f}|_{\mathcal{J}_{\sigma_*}}$ is log-concave (at least
  near the saddle). Apply BL for $\sigma$-concentration:
  $\Var(\sigma(x)) \le 1/(\kappa N)$.
  Here Conlon-Dabkowski~\cite{CD24} extends BL to lattice field
  models with convex gradient action, which is the relevant setting.
\item \textbf{Real coupling on the thimble}: On
  $\mathcal{J}_{\sigma_*}$, the constraint $\im f = \text{const}$
  makes the $\ph$-operator $-\Delta + V(\sigma)$ have a \textbf{real}
  potential $V$. The FK bound / Borell sub-Gaussian estimate gives
  exponential decay.
\item \textbf{Mass gap}: Combine BL concentration + FK bound to
  get $m > 0$ for $N > N_0$.
\end{enumerate}

\subsection{The flowed manifold alternative}

Instead of the exact thimble, one can work with the flowed manifold
$\mathcal{M}_T$ for large but finite~$T$. Advantages:
\begin{itemize}
\item No need to identify all critical points.
\item $\mathcal{M}_T$ is diffeomorphic to $\R^{|\La|}$ (no topology
  change).
\item The sign problem is reduced (not eliminated) but the residual
  oscillations are bounded by $e^{-c(T)N}$ for a function $c(T)$
  increasing in~$T$.
\end{itemize}
For a rigorous proof: take $T$ large enough that the residual
oscillations are smaller than the mass gap signal.

%% ============================================================
\section{Related approaches and results}
%% ============================================================

\subsection{Complex Langevin}

Complex Langevin (CL) dynamics extends stochastic quantization to
complex actions by complexifying the field variables and evolving
with a complex drift. CL avoids the sign problem by construction
but can converge to \textbf{wrong} results.

Boguslavski-Hotzy-Muller~\cite{BHM25} showed that CL converges
correctly when there is a \textbf{single dominant compact thimble}.
Multiple thimbles with different phases cause CL to fail. This
connects CL to the thimble picture and suggests that single-thimble
dominance (which holds at large~$N$) implies CL correctness.

Cai-Dong-Kuang~\cite{CDK20} gave the first rigorous mathematical
analysis of CL, proving that localized probability densities do not
exist for general compact groups, explaining CL instability for
gauge theories.

\subsection{Brascamp-Lieb extensions}

Conlon-Dabkowski~\cite{CD24} extended BL to lattice field models
in $d \ge 2$ with uniformly convex gradient action. They proved
charge-charge correlations in the Coulomb dipole gas are close to
Gaussian, going beyond Dimock-Hurd and Conlon-Spencer. The key
observation: the sine-Gordon measure is invariant for a stochastic
dynamics, allowing a dynamical proof of correlation decay.

This is directly relevant to the BL step on the thimble: the
$\sigma$-measure on $\mathcal{J}_{\sigma_*}$ has a convex gradient
action near the saddle, and Conlon-Dabkowski's extension handles
precisely this setting.

\subsection{P($\ph$)$_2$ via Segal axioms}

Lin~\cite{Lin24} constructed the P($\ph$)$_2$ Gibbs state on
infinite-volume periodic surfaces and proved the mass gap using
Segal's axioms and a ``Bayes principle'' for conditional
probabilities. This is the first construction of an interacting QFT
on surfaces of infinite genus with a mass gap.

This provides an alternative approach to the $N = 1$ test case
(P($\ph$)$_2$ mass gap), complementing the HS/thimble strategy
and the classical Glimm-Jaffe-Spencer approach.

\subsection{Constructive RG}

Giuliani-Mastropietro-Rychkov-Scola~\cite{GMRS24} continued the
rigorous constructive RG program, constructing scale-invariant
response functions at a non-trivial fixed point for fermionic
$\psi^4_d$ (using convergent tree expansion). While the model
differs from ours, the convergent expansion techniques are
potentially applicable to the $1/N$ expansion.

%% ============================================================
\section{Open questions}
%% ============================================================

\begin{enumerate}
\item \textbf{Singularity avoidance}: Can we prove the dominant
  thimble $\mathcal{J}_{\sigma_*}$ avoids $\det(-\Delta+2i\sigma z)=0$?
  The zeros are at $\sigma$ such that $-2iz\sigma$ is an eigenvalue
  of $-\Delta$, i.e., $\sigma = ik^2/(2z)$ (purely imaginary,
  distance $\ge m_0^2/(2z)$ from the real saddle). For large~$N$
  ($z = 1/\sqrt{N}$ small), these are far from the saddle.

\item \textbf{Global log-concavity}: Is $\re f$ concave on
  \emph{all} of $\mathcal{J}_{\sigma_*}$, or only near the saddle?
  The $\Tr\log$ term introduces non-convexity far from the saddle.
  For BL: we need convexity at least within the
  $O(1/\sqrt{N})$-width peak, which holds by the Hessian computation.

\item \textbf{Intersection numbers}: Can we compute $n_{\sigma_*}$
  rigorously? For the dominant saddle on the real axis with $f''$
  positive definite, homotopy arguments should give $n_{\sigma_*}=1$.

\item \textbf{N=1 test case}: In $d = 0$ (one site), $f$ is a
  function of one complex variable. The thimble is a curve in~$\C$
  that can be computed explicitly. Does the thimble analysis
  reproduce the known P($\ph$)$_2$ mass gap?

\item \textbf{Comparison with Kupiainen}: Kupiainen's NLSM proof
  uses the constraint trick (real mass for free) and avoids contour
  deformation entirely. Can the thimble approach give a simpler
  proof of the NLSM mass gap? Or a proof of the \textbf{LSM} mass
  gap (which Kupiainen did not prove)?

\item \textbf{Threshold $N_0$}: What is the minimal $N$ for
  single-thimble dominance? Kupiainen's threshold is
  $N > C\,m_0^{-25}$ (very large). Can the thimble approach
  improve this?
\end{enumerate}

%% ============================================================
\section{Summary}
%% ============================================================

The Lefschetz thimble decomposition provides a natural framework
for the steepest descent contour strategy of \texttt{mass-gap-v3}.
The key facts:

\begin{itemize}
\item The ``steepest descent contour'' $\Gamma = \{\im f = 0\}$
  through the saddle is exactly a Lefschetz thimble.
\item On each thimble, the integrand is proportional to a positive
  function --- the sign problem is solved per-thimble.
\item At large~$N$, a single thimble dominates (saddle-point
  sharpness), so the inter-thimble sign problem is also controlled.
\item The Lawrence-Yamauchi obstruction (for gauge theories) does
  not apply to scalar models: the field space $\R^n$ is contractible.
\item On the thimble, the $\sigma$-measure is real, positive, and
  approximately log-concave --- BL and FK bounds apply.
\item The flowed manifold approach (Fukuma) offers a more robust
  alternative that avoids identifying exact thimbles.
\end{itemize}

The main mathematical challenges are: (1) proving singularity
avoidance, (2) establishing log-concavity on the thimble, and
(3) making the FK bound work with the effective real potential.
These are the subjects of ongoing work in \texttt{mass-gap-v3}
and the Lean formalization.

\begin{thebibliography}{99}
\bibitem{BHM25} K.~Boguslavski, P.~Hotzy, D.~M\"uller,
  ``Lefschetz thimble-inspired weight regularizations for complex
  Langevin,'' arXiv:2412.02396 (2024).
\bibitem{CD24} J.~G.~Conlon and J.~Dabkowski,
  ``Extensions of the Brascamp-Lieb inequality and the dipole gas,''
  arXiv:2412.05122 (2024).
\bibitem{CDK20} Z.~Cai, Y.~Dong, Y.~Kuang,
  ``On the validity of complex Langevin method for path integral
  computations,'' arXiv:2007.10198 (2020).
\bibitem{CSU23} Y.~Cao, R.~Santachiara, R.~Usciati,
  ``On the analytical continuation of lattice Liouville theory,''
  arXiv:2301.07454 (2023).
\bibitem{Fukuma20} M.~Fukuma, N.~Matsumoto,
  ``Worldvolume approach to the tempered Lefschetz thimble method,''
  arXiv:2012.08468 (2020).
\bibitem{Fukuma23} M.~Fukuma,
  ``Simplified algorithm for the Worldvolume HMC,''
  arXiv:2311.10663 (2023).
\bibitem{Fukuma25} M.~Fukuma, N.~Namekawa,
  ``Enhancing the ergodicity of Worldvolume HMC,''
  arXiv:2508.02659 (2025).
\bibitem{Fukuma25b} M.~Fukuma, N.~Namekawa,
  ``Worldvolume HMC for the Hubbard model,''
  arXiv:2507.23748 (2025).
\bibitem{Fukuma26} M.~Fukuma,
  ``Worldvolume HMC for lattice gauge theories,''
  arXiv:2603.28629 (2026).
\bibitem{Gaentgen24} M.~Gaentgen et al.,
  ``Reducing the sign problem with simple contour deformations,''
  arXiv:2407.06896 (2024).
\bibitem{GMRS24} A.~Giuliani, V.~Mastropietro, S.~Rychkov, G.~Scola,
  ``Non-trivial fixed point of $\psi^4_d$, II,''
  arXiv:2404.14904 (2024).
\bibitem{Kupiainen80} A.~Kupiainen,
  Comm.~Math.~Phys.~\textbf{73} (1980) 273--294.
\bibitem{Kupiainen80b} A.~Kupiainen,
  Comm.~Math.~Phys.~\textbf{74} (1980) 199--222.
\bibitem{Lin24} Y.-H.~Lin,
  ``The Bayes principle and Segal axioms for P($\ph$)$_2$,''
  arXiv:2403.12804 (2024).
\bibitem{LY24} S.~Lawrence, S.~Yamauchi,
  ``Convex optimization and contour deformations,''
  arXiv:2311.13002 (2023).
\bibitem{Ostmeyer23} J.~Ostmeyer et al.,
  ``Fermionic sign problem minimization by constant contour shifts,''
  arXiv:2307.06785 (2023).
\bibitem{UA24} M.~Ulybyshev, F.~F.~Assaad,
  ``Beyond the instanton gas approach: dominant thimbles for the
  Hubbard model,'' arXiv:2407.09452 (2024).
\end{thebibliography}

\end{document}
